% GENERATED FILE. Edit canonical lessons, then run npm run generate:books.
% canonical-source-hash: fnv1a64:2d895e1950cd25c2
% canonical-lessons: HI-C306-akal, HI-C306-bhukhmari, HI-C306-sarna, HI-C306-mahamari, HI-C306-haiza

\chapter{Famine and Sickness}
\label{ch:hi-famine-and-sickness}
\hlchaptermodality{306}

\begin{chapteropening}
\textbf{By the end of this chapter:} \emph{I can tell how hunger and disease spread.}

Complete the last lesson of chapter 306: I can tell how hunger and disease spread.
\end{chapteropening}

\section[\texorpdfstring{\hi{अकाल} (akāl)}{अकाल (akāl)}]{\hi{अकाल} (akāl) — a famine}

\label{lesson:HI-C306-akal}



\begin{quote}
Before the new one: say the Hindi for a deluge, then the Hindi for a whirlpool.
\end{quote}

\subsection*{You'll want to know: \hi{अकाल}}
\textbf{\hi{अकाल}} — \emph{akāl} — \textquotedblleft{}a famine\textquotedblright{}.

\begin{quote}
\hi{अकाल} \hi{की} \hi{वजह} \hi{यह} \hi{थी} \hi{कि} \hi{उस} \hi{साल} \hi{बारिश} \hi{नहीं} \hi{हुई।}
\end{quote}

\begin{quote}
\emph{akāl kī vajah yah thī ki us sāl bāriś nahīṃ huī.}
\end{quote}

\begin{quote}
The reason for the famine was that the rains failed that year.
\end{quote}

Masculine.

\begin{cousinweb}
\textbf{\hi{अकाल}} (\emph{akāl}) is Sanskrit \emph{akāla}, \textquotedblleft{}a wrong or bad time\textquotedblright{} (Monier-Williams). In Hindi it means a famine; the dictionary gives the older sense, and the step to a time of hunger is the language's own. A famine \textbf{\hi{पड़ता} \hi{है}}: it falls.
\end{cousinweb}

\begin{grammarlens}[title={What you've built}]
One more word for telling what happened.
\end{grammarlens}

\subsection*{Your turn}
\begin{itemize}
\raggedright
  \item \emph{Say it:} \emph{akāl}
  \item \emph{Say it:} \emph{akāl}, once more
  \item \emph{Say it:} \emph{akāl kī vajah yah thī ki us sāl bāriś nahīṃ huī.} — The reason for the famine was that the rains failed that year.
  \item \emph{Recall:} say the Hindi for a deluge, then the Hindi for a whirlpool, then say \emph{akāl} again
\end{itemize}

\begin{tcolorbox}[breakable,colback=teal!4,colframe=teal!35!black,title={Before you move on}]
What does \hi{अकाल} mean? (\textquotedblleft{}A famine\textquotedblright{}.) Say it once more.
\end{tcolorbox}
\section[\texorpdfstring{\hi{भुखमरी} (bhukhmarī)}{भुखमरी (bhukhmarī)}]{\hi{भुखमरी} (bhukhmarī) — starvation}

\label{lesson:HI-C306-bhukhmari}



\begin{quote}
Before the new one: say the Hindi for a whirlpool, then the Hindi for a famine.
\end{quote}

\subsection*{You'll want to know: \hi{भुखमरी}}
\textbf{\hi{भुखमरी}} — \emph{bhukhmarī} — \textquotedblleft{}starvation\textquotedblright{}.

\begin{quote}
\hi{अकाल} \hi{में} \hi{बहुत} \hi{लोग} \hi{भुखमरी} \hi{से} \hi{मरे।}
\end{quote}

\begin{quote}
\emph{akāl meṁ bahut log bhukhmarī se mare.}
\end{quote}

\begin{quote}
Many people died of starvation in the famine.
\end{quote}

Feminine.

\begin{cousinweb}
\textbf{\hi{भुखमरी}} (\emph{bhukhmarī}) joins \textbf{\hi{भूख}}, hunger, from Sanskrit \emph{bubhukṣā} (CDIAL 9286), with \textbf{\hi{मरना}}, to die: dying of hunger. The vowel of \textbf{\hi{भूख}} shortens inside the word.
\end{cousinweb}

\begin{grammarlens}[title={What you've built}]
One more word for telling what happened.
\end{grammarlens}

\subsection*{Your turn}
\begin{itemize}
\raggedright
  \item \emph{Say it:} \emph{bhukhmarī}
  \item \emph{Say it:} \emph{bhukhmarī}, once more
  \item \emph{Say it:} \emph{akāl meṁ bahut log bhukhmarī se mare.} — Many people died of starvation in the famine.
  \item \emph{Recall:} say the Hindi for a whirlpool, then the Hindi for a famine, then say \emph{bhukhmarī} again
\end{itemize}

\begin{tcolorbox}[breakable,colback=teal!4,colframe=teal!35!black,title={Before you move on}]
What does \hi{भुखमरी} mean? (\textquotedblleft{}Starvation\textquotedblright{}.) Say it once more.
\end{tcolorbox}
\section[\texorpdfstring{\hi{सड़ना} (saṛnā)}{सड़ना (saṛnā)}]{\hi{सड़ना} (saṛnā) — to rot}

\label{lesson:HI-C306-sarna}



\begin{quote}
Before the new one: say the Hindi for a famine, then the Hindi for starvation.
\end{quote}

\subsection*{You'll want to know: \hi{सड़ना}}
\textbf{\hi{सड़ना}} — \emph{saṛnā} — \textquotedblleft{}to rot\textquotedblright{}.

\begin{quote}
\hi{बाढ़} \hi{के} \hi{बाद} \hi{खेतों} \hi{में} \hi{फ़सल} \hi{सड़} \hi{गई।}
\end{quote}

\begin{quote}
\emph{bāṛh ke bād khetoṁ meṁ fasal saṛ gaī.}
\end{quote}

\begin{quote}
After the flood the crop rotted in the fields.
\end{quote}

Past (it …): \hi{वह} \hi{फल} \hi{सड़} \hi{गया}, \emph{vah phal saṛ gayā}.

\begin{cousinweb}
\textbf{\hi{सड़ना}} (\emph{saṛnā}) is set by CDIAL under Sanskrit \emph{śaṭati} (CDIAL 12268), with Hindi \emph{saṛnā}, to rot. Fruit, grain and water \textbf{\hi{सड़ते} \hi{हैं}}.
\end{cousinweb}

\begin{grammarlens}[title={What you've built}]
One more word for telling what happened.
\end{grammarlens}

\subsection*{Your turn}
\begin{itemize}
\raggedright
  \item \emph{Say it:} \emph{saṛnā}
  \item \emph{Say it:} \emph{saṛnā}, once more
  \item \emph{Say it:} \emph{bāṛh ke bād khetoṁ meṁ fasal saṛ gaī.} — After the flood the crop rotted in the fields.
  \item \emph{Recall:} say the Hindi for a famine, then the Hindi for starvation, then say \emph{saṛnā} again
\end{itemize}

\begin{tcolorbox}[breakable,colback=teal!4,colframe=teal!35!black,title={Before you move on}]
What does \hi{सड़ना} mean? (\textquotedblleft{}To rot\textquotedblright{}.) Say it once more.
\end{tcolorbox}
\section[\texorpdfstring{\hi{महामारी} (mahāmārī)}{महामारी (mahāmārī)}]{\hi{महामारी} (mahāmārī) — an epidemic}

\label{lesson:HI-C306-mahamari}



\begin{quote}
Before the new one: say the Hindi for starvation, then the Hindi for to rot.
\end{quote}

\subsection*{You'll want to know: \hi{महामारी}}
\textbf{\hi{महामारी}} — \emph{mahāmārī} — \textquotedblleft{}an epidemic\textquotedblright{}.

\begin{quote}
\hi{उस} \hi{साल} \hi{शहर} \hi{में} \hi{महामारी} \hi{फैली।}
\end{quote}

\begin{quote}
\emph{us sāl śahar meṁ mahāmārī phailī.}
\end{quote}

\begin{quote}
That year an epidemic spread through the city.
\end{quote}

Feminine.

\begin{cousinweb}
\textbf{\hi{महामारी}} (\emph{mahāmārī}) is Sanskrit \emph{mahāmārī} (Monier-Williams), \emph{mahā}, great, and \emph{mārī}, pestilence, a word of the family of \emph{mṛ}, to die. It is a disease that spreads to a whole people.
\end{cousinweb}

\begin{grammarlens}[title={What you've built}]
One more word for telling what happened.
\end{grammarlens}

\subsection*{Your turn}
\begin{itemize}
\raggedright
  \item \emph{Say it:} \emph{mahāmārī}
  \item \emph{Say it:} \emph{mahāmārī}, once more
  \item \emph{Say it:} \emph{us sāl śahar meṁ mahāmārī phailī.} — That year an epidemic spread through the city.
  \item \emph{Recall:} say the Hindi for starvation, then the Hindi for to rot, then say \emph{mahāmārī} again
\end{itemize}

\begin{tcolorbox}[breakable,colback=teal!4,colframe=teal!35!black,title={Before you move on}]
What does \hi{महामारी} mean? (\textquotedblleft{}An epidemic\textquotedblright{}.) Say it once more.
\end{tcolorbox}
\section[\texorpdfstring{\hi{हैज़ा} (haizā)}{हैज़ा (haizā)}]{\hi{हैज़ा} (haizā) — cholera}

\label{lesson:HI-C306-haiza}



\begin{quote}
Before the new one: say the Hindi for to rot, then the Hindi for an epidemic.
\end{quote}

\subsection*{You'll want to know: \hi{हैज़ा}}
\textbf{\hi{हैज़ा}} — \emph{haizā} — \textquotedblleft{}cholera\textquotedblright{}.

\begin{quote}
\hi{बाढ़} \hi{के} \hi{बाद} \hi{गाँव} \hi{में} \hi{हैज़ा} \hi{फैल} \hi{गया।}
\end{quote}

\begin{quote}
\emph{bāṛh ke bād gā̃v meṁ haizā phail gayā.}
\end{quote}

\begin{quote}
After the flood cholera broke out in the village.
\end{quote}

Masculine.

\begin{cousinweb}
\textbf{\hi{हैज़ा}} (\emph{haizā}) is usually traced, through Persian, to Arabic \emph{hayḍa}, \textquotedblleft{}cholera\textquotedblright{}. Bad water after a flood was its usual cause.
\end{cousinweb}

\begin{grammarlens}[title={What you've built}]
One more word for telling what happened.
\end{grammarlens}

\subsection*{Your turn}
\begin{itemize}
\raggedright
  \item \emph{Say it:} \emph{haizā}
  \item \emph{Say it:} \emph{haizā}, once more
  \item \emph{Say it:} \emph{bāṛh ke bād gā̃v meṁ haizā phail gayā.} — After the flood cholera broke out in the village.
  \item \emph{Recall:} say the Hindi for to rot, then the Hindi for an epidemic, then say \emph{haizā} again
\end{itemize}

\begin{tcolorbox}[breakable,colback=teal!4,colframe=teal!35!black,title={Before you move on}]
What does \hi{हैज़ा} mean? (\textquotedblleft{}Cholera\textquotedblright{}.) Say it once more.
\end{tcolorbox}
